\chapter{Architecture Decisions}
\label{ch:decisions}

\newthought{This chapter documents} significant architectural decisions using the Architecture Decision Record (ADR) format. Each decision includes context, rationale, and consequences.

\section{ADR-001: Tauri as Desktop Framework}

\subsection{Context}
OutOf5K requires a desktop application framework that supports:
\begin{itemize}
    \item Local file system access for large demo files
    \item Modern, responsive UI
    \item Reasonable bundle size for distribution
    \item Cross-platform support
\end{itemize}

\subsection{Decision}
Use \textbf{Tauri v2} as the desktop framework.

\subsection{Consequences}

\textbf{Positive:}
\begin{itemize}
    \item Small bundle size (10--20MB vs Electron's 150MB+)
    \item Better performance with system webview
    \item Security-first permission model
    \item Rust backend for performance-critical operations
\end{itemize}

\textbf{Negative:}
\begin{itemize}
    \item Younger ecosystem than Electron
    \item Webview behavior varies across platforms
    \item Rust learning curve
\end{itemize}

\subsection{Alternatives Rejected}

\begin{description}
    \item[Electron] Large bundle size, higher memory usage
    \item[Qt/PyQt] Dated look, smaller ecosystem
    \item[Flutter] Dart unfamiliarity, less mature desktop support
\end{description}

\section{ADR-002: Svelte as Frontend Framework}

\subsection{Context}
Need a frontend framework for complex UI including data visualizations, interactive maps, and responsive design.

\subsection{Decision}
Use \textbf{Svelte 5} with TypeScript.

\subsection{Consequences}

\textbf{Positive:}
\begin{itemize}
    \item Minimal boilerplate compared to React
    \item Excellent performance (compiles to vanilla JS)
    \item Built-in stores, transitions, animations
    \item Strong Tauri integration
\end{itemize}

\textbf{Negative:}
\begin{itemize}
    \item Smaller ecosystem than React
    \item Smaller developer talent pool
    \item Svelte 5 is relatively new
\end{itemize}

\section{ADR-003: Python Subprocess for Demo Parsing}

\subsection{Context}
CS2 demo parsing requires specialized libraries. The best option is \texttt{demoparser2}, a Python library.

\subsection{Decision}
Use Python as a \textbf{subprocess} communicating via stdin/stdout JSON.

\subsection{Consequences}

\textbf{Positive:}
\begin{itemize}
    \item Access to best demo parser (demoparser2)
    \item Python ecosystem for future ML features
    \item Process isolation (parser crashes don't affect main app)
    \item Independent development and testing
\end{itemize}

\textbf{Negative:}
\begin{itemize}
    \item Process spawn overhead
    \item Distribution complexity (bundle Python or require installation)
    \item IPC latency
\end{itemize}

\subsection{Distribution Strategy}
Use PyInstaller to create a standalone Python executable bundled with the app.

\section{ADR-004: SQLite as Primary Database}

\subsection{Context}
Need persistent storage for demos, events, statistics, and settings. Desktop app requires zero-configuration setup.

\subsection{Decision}
Use \textbf{SQLite} accessed via sqlx in Rust.

\subsection{Consequences}

\textbf{Positive:}
\begin{itemize}
    \item Zero configuration required
    \item Single file (easy backup/migration)
    \item Excellent performance for single-user workloads
    \item Full SQL support including JSON functions
\end{itemize}

\textbf{Negative:}
\begin{itemize}
    \item Single-writer limitation
    \item No built-in replication for team features
    \item Limited full-text search
\end{itemize}

\section{ADR-005: PlantUML for Diagrams}

\subsection{Context}
Documentation requires architecture and UML diagrams that are version-controlled and LaTeX-compatible.

\subsection{Decision}
Use \textbf{PlantUML} for all technical diagrams.

\subsection{Consequences}

\textbf{Positive:}
\begin{itemize}
    \item Text-based, version controllable
    \item Excellent C4 and UML support
    \item LaTeX-compatible outputs (PDF, PNG, SVG)
    \item Wide tool support (VS Code, IntelliJ, CLI)
\end{itemize}

\textbf{Negative:}
\begin{itemize}
    \item Requires Java
    \item Learning curve for syntax
    \item No native GitHub rendering (unlike Mermaid)
\end{itemize}

\section{ADR-006: Phased 3D Implementation}

\subsection{Context}
Users want 3D replay capability, but building a full 3D viewer is complex and time-consuming.

\subsection{Decision}
Implement visualization in phases:
\begin{enumerate}
    \item Phase 1 (MVP): 2D radar-style map
    \item Phase 2: Enhanced 2D with paths, utility, heatmaps
    \item Phase 4: 3D replay viewer (scope TBD)
\end{enumerate}

\subsection{Consequences}

\textbf{Positive:}
\begin{itemize}
    \item Ship faster with 2D MVP
    \item Validate core features before 3D investment
    \item User feedback informs 3D priorities
    \item 2D is still valuable for analysis
\end{itemize}

\textbf{Negative:}
\begin{itemize}
    \item Not full-featured initially
    \item Less exciting for marketing
\end{itemize}

\subsection{Success Criteria}
Proceed to 3D only when:
\begin{itemize}
    \item Users explicitly request 3D features
    \item Technical feasibility is proven
    \item Resources are available
\end{itemize}
